\begin{CJK}{UTF8}{gbsn}
\chapter{摘要}
%\addcontentsline{toc}{chapter}{Introduction}  
%本论文通过希格斯玻色子衰变到四轻子过程研究了希格斯粒子的性质。基于CMS探测器收集的质心系能量为13 TeV的质子质子对撞数据，测量了希格斯至四轻子衰变的基准截面，包括积分截面和微分截面。这些测量使用了CMS第二期运行获取的全部数据，即质心系能量为13 TeV、积分亮度为137.1 $\ifb$的数据，在测量中希格斯质量取值为125.09 $\GeV$。测量结果直接检验了标准模型希格斯相关的在TeV能标上的QCD微扰计算。为了使所测量结果不依赖于希格斯的理论预测，论文基于末态轻子运动学信息和事例拓扑结构定义了基准相空间，所有的截面测量都限定与基准相空间。在质心系能量为13 TeV情况下，积分基准截面测量值为$2.73^{+0.30}_{-0.29}=2.73^{+0.23}_{-0.22}({\rm stat.})^{+0.24}_{-0.19}({\rm syst.})~{\rm fb}$。 另外针对多个对希格斯玻色子产生机制敏感的运动学观测量，包括快度、四轻子体系的横动量、相关的多重喷注以及领头喷注的横动量等进行了基准微分截面测量。测量结果对于检验标准模型希格斯粒子的预测至关重要，同时可以用于约束许多超标准模型理论的相空间。在误差范围内，积分截面和微分截面结果均与标准模型理论预言一致。

基于CMS探测器收集的 质心系能量 =13TeV, 质子-质子对撞数据，在$ H \to 4\ell$ decays ($\ell = \Pe$, $\mu$) 衰变道 测量Higgs粒子的基准截面 (积分和微分)。测量中使用Higgs质量 以及大型强子对撞机（LHC）产生的积分亮度为 （Full Run II）的实验数据。
此次测量是为精确检验粒子物理标准模型（SM），经过长时间的研究，物质及物质间的相互作用被标准模型精确描述（电磁力，弱力，强力不包括引力）。标准模型提供的数学描述只能适用于质量为0 的粒子。Higgs机制的引入，使得标准模型更加普世，解决了标准模型只能描述0质量粒子的问题。在Higgs机制中，粒子通过电弱对称性破缺获得质量。Higgs粒子的存在就是Higgs机制的直接证明。
在Tevatron，LEP高能物理实验的精确测量中，已经证明标准模型的正确性，但是仍然没有观测到Higgs粒子。在2010年，大型强子对撞机被建造，用来对标准模型做更多的检验，以及验证之前的实验结果。LHC安装在CERN（欧洲核子研究中心）坐落在瑞士-法国边境，日内瓦近郊。它是在粒子物理同类实验中，能量最高的实验装置。它有足够的能力将质子束流加速到14TeV 亮度峰值为。它是人类有史以来建造过的最先进，最复杂和最大的实验装置，由来自超过100个国家，10000名工程技术人员以及科学家共同设计建造。从1998年到2008年，LHC被建造在地下75-150米，3.8米宽，27千米长的隧道里，该隧道在1983-1988年曾经用来安装LEP对撞机。在LHC中两个反向旋转的质子束流运行在高真空的管道中以避免束流和空气分子的额外碰撞。2012年，在LHC上运行的两个实验CMS和ATLAS联合发布发现了一个新的重标量玻色子，Higgs玻色子。基于这个发现，物理学家希望通过精确测量Higgs粒子的性质，以进一步检验标准模型。Higgs粒子的质量是Higgs机制中的唯一自由参数，它的测量结果是 。
特定过程的产生截面是Higgs粒子的重要性质。它给出该过程发生的机率。这篇论文，通过CMS收集的　质心系能量 =13TeV, 积分亮度为 　的　质子-质子对撞数据，在　$ H \to 4\ell$ decays ($\ell = \Pe$, $\mu$) 衰变道测量Higgs粒子的积分，微分基准产生截面．考虑到以下因素Included　Higgs微分截面测量基于Higgs运动学，Higgs横向动量可以用于检测ｐＱＣＤ的计算（特别是当实验中使用全举多ｊｅｔ），另外这个变量对Higgs的拉适量的结构非常敏感。Higgs的快度对该实验中的质子部分子分布函数（PDFs）非常敏感。jet相关的观测量对理论模型非常敏感，boosted夸克，胶子辐射也在这篇论文中讨论。目前的测量结果与标准模型预言相一致，但是，在此类测量结果中，任何与标准模型的偏差，都直接指向超越标准模型的贡献。这些测量结果也是对微扰QCD计算的在TeV能标的直接检验。
	通过极大似然法拟合参数化的四轻子不变质量谱来测量在 $pp\to$H$\to4\ell$ 过程中的　积分和微分基准截面，其中 $N_{\mathrm{obs}}(m_{4\ell})$ 和基准截面 ($\sigma_{\mathrm{fid}}$) 在拟合过程中直接求得．系统误差通过整合到似然函数中的冗余参数表示．结果通过 asymptotic和基于profile　likelihood　ratio的统计检验获得．从观测到的分布中求解探测器效应的过程与 Ref[14-15].相同．在微分截面的测量中，探测器效率和探测器分变率由探测器响应矩阵描述，这个矩阵描述了基准水平的事例和重建水平事例之间的差别．这个矩阵由对角线矩阵元主导，适当的非对角元反应了ｊｅｔ相关的观测量．
实验不确定性中的轻子识别效率和亮度测量精度是该分析中最重要的不确定性，理论不确定性并不占主导作用．为了评价该测量对理论模型依赖程度，使用不同的探测器相应矩阵来重复求解探测器效应的过程，这些矩阵通过在实验限制的范围内改变标准模型不同产生模式之间的比例来获得．最终发现该不确定性相比于实验不确定性是可以忽略不记得．
	为了保证所有得测量相互独立而不依赖于Higgs得产生方式，所有得截面测量在一个被轻子动力学和事例拓扑结构定义得基准相空间进行．积分基准截面的测量结果是$2.73^{+0.30}_{-0.29}=2.73^{+0.23}_{-0.22}({\rm stat.})^{+0.24}_{-0.19}({\rm syst.})~{\rm fb}$．微分基准截面的测量结果依赖于Higgs的产生方式，其结果是不同动力学变量的函数，比如：四轻子末态的赝快度，横动量，ｊｅｔ多样性以及领头阶ｊｅｔ横动量．微分散射截面在观测量的不同区间进行．每一个动力学变量或者微分观测量的区间边界的确定都被优化，以使得在观测量的每一个区间有相近的期望测量值．在微分散射截面的测量中，观测量的每一个区间中的背景事例被减去．背景事例的质量谱与所有微分区间质量谱都不相同，因为这里没有共振态结构．在微分散射截面的测量中，对于irreducible和reducible背景，通过蒙特卡洛样本和control　region的实验数据分别获得 Ref[15]．该测量对Higgs机制的检验非常重要．同时也给超出标准模型的理论提出限制．
所有的积分，微分测量结果与标准模型理论预言在误差范围内相一致．需要更精确的Higgs性质测量，任何与标准模型的偏差都是新物理的证据．

%	
%	
%	
%	
%	
%	
%	该论文研究了在希格斯玻色子四轻子衰变道中其性质的测量。在质心系能量为13TeV的质子质子对撞机中，使用CMS探测器收集的数据，测量希格斯至四轻子衰变的基准截面，其包括积分截面和微分截面。测量中固定希格斯质量为125.09GeV，使用质心系能量13TeV积分亮度为137.1$\ifb$的数据（第二期运行阶段全部数据）。测量结果直接检验了在TeV能级上标准模型希格斯相关的QCD微扰计算。
%为了使所有测量结果独立于希格斯的产生方式，基于轻子运动学和事件拓扑结构定义基准相空间，所有的截面测量都取自基准相空间。在质心系能量13TeV下，测量的积分基准截面为$2.73^{+0.30}_{-0.29}=2.73^{+0.23}_{-0.22}({\rm stat.})^{+0.24}_{-0.19}({\rm syst.})~{\rm fb}$.
%基准微分截面考察了多个对希格斯玻色子产生机制敏感的运动学观测量，包括快度，四轻子体系的横动量，相关的多重喷注以及领头喷注的横动量。测量结果对于检验标准模型希格斯部分的预测至关重要，并且可以用于许多超标准模型理论的相空间约束。
%在误差范围内，积分截面和微分截面结果均与标准模型理论预言一致。
%	
%	
%	
%	
%	
%	
%	
%	
%	
%	
%	
%	
%	
%	
%	
%	
%	该论文研究了在质子-质子对撞中，质心系能量分别为7,8以及13TeV下，通过四轻子衰变 $ H \to 4\ell$ ($\ell = \Pe$, $\mu$) 的希格斯玻色子产生截面的测量问题。
%希格斯玻色子的基准产生截面测量数据都源自于大型强子对撞机LHC上运行的CMS探测器，对应于不同质心系能量，积分亮度分别为7TeV下5.1~\fbinv, 8TeV下19.7~\fbinv, 以及13TeV下
%2.8~\fbinv。这些测量结果为标准模型在TeV能级下的希格斯物理相关的QCD微扰计算提供了突破口。测量结果与标准模型的理论预期有任何的偏离都将会为超出标准模型的新物理提供直接的线索。
%\par 
%质心系能量为8TeV时的基准微分截面已给出，对希格斯玻色子产生机制敏感的几个动力学观测量进行了考察，包括快度，四轻子体系的横动量，相关的多重喷注，领头喷注的横动量，以及对领头喷注和希格斯玻色
%子候选粒子间的分离。对希格斯玻色子衰变机制敏感的几个动力学观测量的研究结果也已公布~\cite{Reference40}，这些动力学观测量包括：两组轻子对的不变质量，以及Collins-Soper体系下轻子形成的五个螺旋度角。在希
%格斯的测量基础上，本文也利用了同样的质心系能量为8TeV的数据和微分变量，来对 $Z \rightarrow 4\ell$ 的基准微分截面进行测量。 $Z \rightarrow 4\ell$的基准微分截面的测量对检验标准模型的
%理论预言以及对希格斯玻色子产生截面的测量方法的合理性的验证都至关重要。
%\par $ H \to 4\ell$ 和 $Z \rightarrow 4\ell$ 的积分基准截面比率的计算结果也已发表。所有的截面测量都取自于基准相空间，其定义是来自于轻子运动学和事例拓扑结构的筛选的结果。总的基
%准截面测量结果在7, 8和13TeV下分别为$0.56^{\rm +0.67}_{\rm -0.44}({\rm stat.})^{\rm +0.21}_{\rm -0.06}({\rm sys.})^{+0.02}_{-0.02}({\rm model})$ ${\rm fb}$,
%$1.11^{\rm +0.41}_{\rm -0.35}({\rm stat.)^{\rm +0.14}_{\rm -0.10}({\rm sys.}) ^{+0.08}_{-0.02}({\rm model}){\rm fb}$
%以及 $2.48^{+1.45}_{-1.12}({\rm stat.})^{+0.28}_{-0.17}({\rm sys.})^{+0.01}_{-0.04}({\rm model})~{\rm fb}$。所有的截面测量结果和基于标准模型是理论预期之间没有显著性的偏差。
\end{CJK}
